\subsection{Five continuous payments for a retained reference square}
\label{sec:retained-reference-payments}
\begin{lemma}\label{lem:retained-reference-payments}
For the functions $\mathcal U_i$ below, one has
$\mathcal U_i(d)\le1-\varepsilon_i d$ throughout $0\le d\le D_i$.
\begin{align*}
 \mathcal U_1(d)={}&e^{-37d/4}
   +\frac{39}{4}d e^{-4d-2d^2}
   -\frac34d e^{-851d/96}
   -9d^2(\kappa-2d)^2e^{-7d/3},\\
 \mathcal U_2(d)={}&e^{-159d/16}
   +\frac{261}{26}d e^{-9d/2-81d^2/32}
   -\frac34d e^{-10d},\\
 \mathcal U_3(d)={}&e^{-203d/16}
   +\frac{72}{7}d e^{-15d/4-81d^2/32},\\
 \mathcal U_4(d)={}&e^{-209d/16}
   +9d e^{-25d/8-2d^2}
   +(d/4+4d^2)e^{-349d/32},\\
 \mathcal U_5(d)={}&e^{-25d/2}
   +\frac{91}{10}d e^{-7d/2}.
\end{align*}
Here $\kappa=19/32$, and the parameters are
\[
 \begin{array}{c|cc|cc}
  i&D_i&\varepsilon_i&\deg P_i&\min_j b_{i,j}\text{ lower bound}\\\hline
  1&1/4&1/800&32&1/100\\
  2&1/5&1/16&32&1/5\\
  3&1/5&1/160&32&1/10\\
  4&1/4&1/600&32&1/20\\
  5&1/4&1/60&16&1/100.
 \end{array}
\]
The polynomials $P_i$ and their Bernstein coefficients $b_{i,j}$ are
defined in the proof, on one interval for each row.
\end{lemma}
\begin{proof}
Put
\[
 E(z)=\sum_{j=0}^{16}\frac{(-z)^j}{j!},\qquad
 O(z)=\sum_{j=0}^{15}\frac{(-z)^j}{j!}.
\]
The signed Taylor remainder gives $O(z)\le e^{-z}\le E(z)$ for
every $z\ge0$. In the displayed $\mathcal U_i$, replace an exponential
with a positive prefactor by $E$, and one with a negative prefactor by
$O$; call the resulting polynomial $\widehat{\mathcal U}_i$.
All prefactors have the indicated sign on their stated intervals.
The polynomial
\[
 P_i(d)=\frac{1-\varepsilon_i d-\widehat{\mathcal U}_i(d)}{d}
\]
has a removable zero in the numerator at $d=0$.
If $P_i(d)=\sum_{k=0}^{n_i}a_{i,k}d^k$, its coefficients in
the Bernstein basis on $[0,D_i]$ are exactly
\[
 b_{i,j}=\sum_{k=0}^{j}
  a_{i,k}D_i^k\frac{\binom jk}{\binom{n_i}k},\qquad 0\le j\le n_i.
\]
Substitution of the five displayed formulas gives the 149 positive
coefficients with the rational lower bounds in the table.
 Applying
Lemma~\ref{lem:bernstein-certificate} gives $P_i\ge0$ on the entire
interval. Therefore
$\mathcal U_i\le\widehat{\mathcal U}_i\le1-\varepsilon_i d$,
including the endpoint $d=0$.
\end{proof}
